% Corrigé intégré automatiquement pour la banque Exercices.
% Source : B2_ProposerModele/B2_16_Hyperstatisme/71_Robovolc_02/71_Robovolc_02.tex
\begin{corrigebox}[Corrigé]
\CorrigeQuestion{1}
On identifie les classes d'équivalence cinématique, les surfaces de contact et les mouvements relatifs autorisés. Chaque contact est remplacé par la liaison normalisée correspondante, puis les axes et points caractéristiques sont reportés sur le schéma cinématique minimal.
\CorrigeQuestion{2}
On recense les inconnues de liaison réellement indépendantes et le nombre d'équations d'équilibre disponibles après prise en compte des mobilités. Le degré d'hyperstatisme est obtenu par le bilan $h=N_s-r_s$, puis contrôlé en identifiant les inconnues redondantes.
\CorrigeQuestion{3}
On isole le solide ou l'ensemble indiqué, on dresse le bilan complet des actions mécaniques extérieures et on choisit un point de réduction qui élimine le maximum d'inconnues. Le PFS donne $\sum\vec F=\vec0$ et $\sum\vec M_A=\vec0$ ; les projections utiles fournissent les inconnues, puis leur signe permet de vérifier le sens supposé.
\end{corrigebox}
